% Recommended paper title:
% Sparsity Spillover in Transformers: Local Pressure, Model-Wide Consequences

\section{Introduction}

Activation sparsity offers a route to reducing the cost of Transformer inference~\cite{liu2025teal,zhang2025rsparse}. For a linear transformation \(y=xW\), any zero activation coordinate \(x_j=0\) nullifies the corresponding products and can, in principle, be skipped by a specialized kernel~\cite{lee2024cats,liu2025teal,zhang2025rsparse}. Unlike weight sparsity, which is fixed after training, activation sparsity is input-dependent and can vary across tokens, layers, and model components. Recent training-free methods such as TEAL and R-Sparse showed that pretrained Transformers can tolerate thresholding sufficiently small activations to zero with limited degradation in model quality~\cite{liu2025teal,zhang2025rsparse}. These methods therefore establish a post-training quality--compute frontier: increasing the threshold makes a larger fraction of matrix-multiplication operations removable, but also degrades model quality~\cite{lee2024cats,he2024chess,liu2025teal,zhang2025rsparse}. We ask whether interventions applied during pretraining can instead shape activation distributions as the model learns and extend this frontier toward lower computation at the same model quality.

A direct way to reshape activation distributions is to apply an \(\ell_1\) penalty to selected activations, pressuring them toward zero. Existing work on this intervention has focused on activation sites within Transformer feed-forward networks (FFNs). ProSparse applies it during post-training adaptation: it replaces the relevant activations with ReLU and progressively strengthens an \(\ell_1\) penalty on the intermediate activations that feed each FFN's output projection~\cite{song2025prosparse}. More recently, Cetin et al. showed that the simpler combination of ReLU and \(\ell_1\) pressure at the same FFN sites can induce over \(99\%\) sparsity with negligible quality degradation~\cite{cetin2026sparser}. The same work released specialized sparse GPU kernels and demonstrated measured wall-clock speedups in both inference and training, establishing that the induced FFN sparsity yields runtime gains. However, the reported sparsity in both works characterizes only the targeted FFN activations; it is neither a global sparsity measurement nor a model-wide estimate of removable computation. Neither work measures how untargeted components, including attention, respond to the local pressure. This leaves both a behavioral and a computational gap: \emph{when one set of activations is pushed toward zero, does the rest of the network remain unchanged, and what is the resulting effect on the entire model?}

We study this question during pretraining in Pythia, a family of Transformer language models trained under a common setup and released at several sizes~\cite{biderman2023pythia}. Within this setting, we apply an \(\ell_1\) penalty to FFN activations, vary its strength \(\lambda\), and replace GeLU with ReLU. Locally, the intervention works: stronger pressure moves the targeted FFN activations toward zero. The response, however, is not confined to the FFNs: as those activations concentrate near zero, the unregularized attention activation distributions broaden. We call this opposing, nonlocal response \emph{sparsity spillover}. Spillover matters computationally because local FFN statistics alone do not determine how much matrix-multiplication work becomes removable across the model's entire forward pass.

To quantify the implications for inference compute, we define \(\Rmodel{z}\), the percentage of matrix-multiplication operations across the model that can be avoided under an exact-zero activation pattern \(z\). We call each activation tensor that serves as an input to a matrix multiplication and is eligible for sparsification a \emph{sparsification site}. \(\Rmodel{z}\) therefore converts observed zeros at the selected sites into a model-wide operation count rather than averaging local sparsity percentages. For a specified set of sparsification sites and a fixed sequence length \(T\), we further define \(\Rmodelmax\) as the corresponding architecture-specific ceiling: the percentage of model-wide matrix-multiplication operations that would be avoided if activations at those sites were fully sparse. Here, \(T\) denotes the number of tokens processed in the forward pass for the compute accounting. Unlike \(\Rmodel{z}\), \(\Rmodelmax\) can be computed a priori once the architecture, candidate sites, and \(T\) are fixed, independently of the input examples and observed activation values. This ceiling depends strongly on both architecture and sequence length. For example, when sparsification is restricted to FFN sites at \(T=2048\), \(\Rmodelmax\) increases from \(11.8\%\) for Pythia-14M to \(86.8\%\) for Pythia-12B~\cite{biderman2023pythia}. For the same Pythia-12B model, increasing the sequence length from \(T=2048\) to \(T=50{,}000\) lowers the ceiling from \(86.8\%\) to \(49.9\%\), as quadratic attention cost grows relative to the FFN projections. Thus, the computational value of a sparsification target depends on where it is placed, the model architecture in which it appears, and the sequence length at which it is evaluated---not only on its local sparsity percentage.

% Empirical scope and provenance: tentative Findings F001/F002, sourced from
% Runs 011--013 and Analyses 007/008 (one seed, Pythia-14M, one MiniPile pass).
Motivated by these observations, we propose an architecture-wide sparsification methodology that operates during pretraining. Rather than treating FFN activations in isolation, the methodology uses \(\Rmodel{z}\) and \(\Rmodelmax\) to evaluate candidate sparsification sites across the network and place interventions where they can affect model-wide computation. It combines \(\ell_1\) regularization at strategic FFN sites with thresholding at selected attention sites, building on prior work that applies magnitude thresholding to hidden states and attention computations~\cite{lee2024cats,he2024chess,liu2025teal}. For the signed attention activations, the method includes two threshold forms: a one-sided threshold, which zeros all negative values and positive values below \(\kappa\), and a symmetric threshold, which zeros values whose magnitude is below \(\kappa\). Gradient orthogonalization is also an integral component of the methodology. We introduce this component because the gradients of the task loss and sparsity objective can point in conflicting directions. Motivated by gradient-surgery methods~\cite{yu2020pcgrad}, we project the sparsity gradient so that it no longer opposes the task gradient. In the current one-seed Pythia-14M full-pass evidence, adding four-site orthogonal \(\ell_1\) pressure to the matched A4-Z threshold topology improves both validation loss and measured logical-product opportunity for \(\kappa\le0.1\); the advantage contracts with threshold strength and reverses in validation loss at \(\kappa=0.5\). Expanding A4-Z with symmetric post-RoPE Q/K/V gates is near-null at \(\kappa=0\), improves both matched axes at \(\kappa=0.01\), and provides increasing logical opportunity with increasing validation-loss cost at larger thresholds. These results are initial evidence for a better model-wide quality--logical-compute frontier, not scale-independent or measured-speedup claims.

In summary, our contributions are:

\begin{itemize}
    \item We identify \textbf{sparsity spillover}: during pretraining, FFN sparsity pressure moves the targeted activations toward zero but is accompanied by fewer near-zero values and broader magnitude distributions in unregularized attention activations. This shows that local FFN statistics do not capture the full model response.

    \item We introduce model-wide sparsity accounting over eligible sparsification sites. \(\Rmodel{z}\) measures the share of the model's matrix-multiplication work made removable by observed exact zeros, while \(\Rmodelmax\) gives an architecture-specific ceiling for a chosen set of sites and sequence length. Because \(\Rmodelmax\) is computable a priori, this accounting supports architecture-aware site selection and exposes how the value of a target depends on its location, model scale, and sequence length.

    \item An architecture-wide sparsification methodology applied during pretraining. It combines \(\ell_1\) regularization at strategic FFN sites, spillover-guided thresholding at selected attention sites with two threshold forms, and gradient orthogonalization to address interference between the task and sparsity objectives. At Pythia-14M, matched one-seed evidence indicates that conflict-aware pressure improves A4-Z validation loss and logical-product opportunity at moderate thresholds, while adding symmetric post-RoPE Q/K/V gates extends the attainable logical-opportunity range with a dose-dependent quality cost. Replication and larger-scale evaluation remain necessary.
\end{itemize}
